\chapter{Einleitung}

In der Welt der digitalen Signalverarbeitung ist die Faltung (engl.: \glqq Convolution\grqq) ein mathematischer Vorgang, bei dem 
ein Eingangssignal mit einem weiteren Signal, einer sog. Impulsantwort gefaltet wird.
Eine Impulsantwort ist hierbei als diskretes Audiosignal zu verstehen, das die Klangeigenschaften
eines bestimmten Systems (bspw. Lautsprecherboxen, Räume, Filter, ... etc.) beinhaltet.
Die Convolution beschreibt also mathematisch den Einfluss eines solchen Systems auf ein Eingangssignal.\cite[Kapitel 5]{lyons2010}

In den folgenen Kapiteln werden folgende Aspekte der Convolution erläutert:
\begin{itemize}
    \item Die mathematischen Grundlagen der Convolution diskreter Audiosignale
    \item Veranschaulichung verschiedener Systeme, die mithilfe der Convolution realisiert werden können
    \item Praktische Anwendung und die damit verbundenen Problemstellungen
    \item Gegenüberstellung: Naive Umsetzung der diskreten Convolution und Umsetzung mithilfe der \textit{Fast Fourier Transformation} (kurz: FFT)
    \item Implementierung in Software zur digitalen Audiobearbeitung
\end{itemize}

Allerdings wird im Rahmen dieser Facharbeit nicht auf die mathematische Definition der
die schnelle Fourier Transformation eingegangen, weil dies den Rahmen der Seminararbeit sprengen würde.
Stattdessen werden alle Techniken und Theoreme erläutert, die für das grundlegende
Verständnis der Convolution notwendig sind.